\begin{center}
    \large
    \begin{singlespace}    
        \textbf{\chineseTitle{}} \\[0.5cm]
    \end{singlespace}

    \begin{singlespace}    
    學生：\studentCnName  \hspace{2.5cm}  指導教授：\advisorCnName \hspace{0.1cm} 博士 \\

    \ifdefined\advisorCnNameB % 如果有共同指導教授
        \hspace{9.6cm}  \advisorCnNameB ~ 博士 \\
    \fi
    \end{singlespace}

    \vspace{0.5cm}

    \begin{singlespace}
        國立陽明交通大學\\
        \NameofDepartmentInstituteCN\\[0.2cm]
    \end{singlespace}

    \textbf{摘\hspace{1cm}要} \\[0.5cm]

\end{center}

\normalsize 

隨著5G與未來6G網路演進，網路切片技術已成為支援垂直產業多樣化服務品質保障之關鍵。然而，傳統跨網域端到端切片編排依賴命令式介面，面臨配置繁瑣與協調遲滯等技術挑戰。為此，本研究遵循3GPP與ETSI之服務化管理架構（SBMA）規範，設計並實作一套端到端意圖驅動網路切片編排系統。系統採用宣告式Kubernetes資源模型（KRM），以自定義資源作為北向意圖管理服務介面，向外部消費者隱蔽底層技術細節，達成以目標為導向的自動化運維範式。

在核心元件實作上，本研究自主開發端到端編排器，實質承擔網路切片管理功能（NSMF）與意圖處理器職能。針對南向異質子網域，系統採取雙軌編排機制：無線存取網網域整合 Nephio 編排平台落實自動化之套件生命週期管理；核心網網域則採用混合式編排機制將轉譯之參數注入核心網中。此外，工作負載叢集之自定義 Operator 則精確扮演子網管理功能（NSSMF）角色，負責本地資源之遞迴調和與實體化配置。

經由端到端實驗驗證，結果證實本系統能精確將高階服務等級協議（SLA）期望自動轉譯為底層無線資源與核心網策略。系統在資源競爭時能有效保障關鍵型切片之數據面服務品質，並透過狀態回饋機制展現優異的最終一致性調和與配置自癒能力，為未來行動通訊網路向零接觸管理（Zero-touch Management）演進提供具實踐價值的開源參考架構。

\vspace{1cm}

% 中文摘要及關鍵詞 5-7 個 
\textbf{關鍵字：}5G、意圖驅動管理、網路切片、Nephio、OCUDU
